\chapter{The repository, file by file}\label{ch:repo}

Everything in the previous fourteen chapters is on one branch of \texttt{bkrobust},
\texttt{ze/experiments}. This chapter is the map: how the branch is built, where each chapter's
files are, what to run first, and the caveats I know about. The same map, in markdown, is
\texttt{ZE\_EXPERIMENTS.md} at the root of the branch.

\section{How the branch is built}

\begin{antes}
\antesq[pre-registro]{Why is the branch built on top of your history rather than as a fresh repository?}
\antesq[quadro-observavel]{Which of the five commits would you read first if you only cared about Table~4?}
\end{antes}

\begin{ideia}
The branch shares your history, so \texttt{git merge} works and \texttt{git diff} against your
last common commit shows exactly what I added. Nothing on it rewrites a file of yours without a
commit that says why.
\end{ideia}

The branch starts from \texttt{iclr2027-eval-plan} (commit \texttt{ba7ddc1}), the evaluation plan
that is already on GitHub. On top of it come six steps, drawn in the diagram below: the merge of the
August branch \texttt{rho-breakdown-evidence}, which I never pushed, then five commits.

\begin{center}
\begin{tikzpicture}[node distance=4mm and 6mm]
\node[blkc, text width=3.3cm] (base) {\texttt{d3cd71f}\\ your \texttt{main} on 11 Sep};
\node[blk, text width=3.3cm, right=of base] (plan) {\texttt{ba7ddc1}\\ evaluation plan (pushed)};
\node[blko, text width=3.3cm, right=of plan] (aug) {merge: August package\\ \texttt{rho\_breakdown/}};
\node[blk, text width=3.3cm, below=8mm of base] (c1) {1\quad \texttt{src/}, \texttt{tests/}};
\node[blk, text width=3.3cm, right=of c1] (c2) {2\quad \texttt{experiments/}, \texttt{docs/}};
\node[blk, text width=3.3cm, right=of c2] (c3) {3\quad \texttt{results/}};
\node[blkg, text width=5.2cm, below=8mm of c1.south west, anchor=north west] (c4) {4\quad \texttt{rho\_breakdown/}: dated folders, 14 Aug to 16 Sep};
\node[blkg, text width=5.2cm, right=6mm of c4] (c5) {5\quad \texttt{ZE\_EXPERIMENTS.md}, this polycopié};
\draw[flow] (base) -- (plan); \draw[flow] (plan) -- (aug);
\draw[flow] (aug.south) |- ++(0,-3mm) -| (c1.north);
\draw[flow] (c1) -- (c2); \draw[flow] (c2) -- (c3);
\draw[flow] (c3.south) |- ++(0,-3mm) -| (c4.north);
\draw[flow] (c4) -- (c5);
\end{tikzpicture}
\end{center}

Commits 1 to 3 are the protocol run of Chapters~\ref{ch:where} to~\ref{ch:tables}, as it sat
uncommitted in my working copy of your repository: 158 files, split into library changes, scripts
with docs, and results. Commit~4 brings in the experiment folders that lived outside the repository.
Commit~5 is the index and this document.

\begin{medido}[The new tests, \texttt{tests/\{benchmarks,criterion,estimation,knowledge\}}]
\texttt{pytest -q} on the nine new test files of the protocol run: \textbf{39 passed} in 4.4\,s,
16 September, with the library as committed on the branch. They cover the audit oracle (back-door
against GAC), the discovery arm, measurement selection, the sortability probe, the closed-form
$\Ostar$, the knowledge interval and the ancestral block.
\end{medido}

That measurement covers the new tests only. Your older tests were not part of it, so run the whole
suite once after merging, before trusting that the library changes left your results untouched.

\begin{noprato}[Getting it]
\texttt{git fetch origin ze/\allowbreak{}experiments}, then \texttt{git diff d3cd71f origin/\allowbreak{}ze/\allowbreak{}experiments --stat}
to see what changed in your package. Your \texttt{main} has moved since \texttt{d3cd71f}, so merge,
do not rebase; the library changes touch \texttt{benchmarks/measure.py}, \texttt{hybrid.py} and
\texttt{search/exact*.py}, which are the places a conflict would appear.
\end{noprato}

\prompt[conceito=pre-registro,dificuldade=0.8]{Which commit on the branch holds the
pre-registrations of the protocol run, and why are they in the same commit as the tables?}{Commit~3,
\texttt{results/}. Each results folder carries the \texttt{PREREGISTRATION*.md} written before its
rows were read, next to the tables built from those rows, so that a reader can check the predictions
against the outcomes in one place (for example \texttt{results/\allowbreak{}ledger/\allowbreak{}PREREGISTRATION.md} with
\texttt{TABLE4.md}).}

\section{Where each chapter lives}

\begin{antes}
\antesq[tabelas-pagas]{Which folder holds the script that reproduces the 27 cells of panel~A?}
\antesq[gac-backdoor]{Where is the test that stops an audit from using the back-door criterion?}
\end{antes}

The table below is the index I would want if I were you. Paths are relative to the repository root.

\begin{center}
\small
\begin{tabular}{@{}p{1.6cm}p{6.1cm}p{6.6cm}@{}}
\toprule
chapter & folder & start with \\
\midrule
\ref{ch:pilot} & \texttt{rho\_breakdown/pilot/}, \texttt{e1prime/} & \texttt{REPORT.md}, \texttt{NUMBERS.md}, \texttt{e1prime/RESULTS.md} \\
\ref{ch:locality} & \texttt{rho\_breakdown/\allowbreak{}locality/}, \texttt{x3\_skeleton/}, \texttt{rerank\_2026-08-14/} & \texttt{locality/VERDICT.md}, \texttt{x2/\allowbreak{}cfirewall/\allowbreak{}STATUS.md}, \texttt{x3\_skeleton/\allowbreak{}RESULTS.md} \\
\ref{ch:redefinition} & \texttt{rho\_breakdown/\allowbreak{}radius\_redef\_2026-09-07/}, \texttt{rnull\_real\_2026-09-09/} & \texttt{EXPERIMENTS.md}, \texttt{IMPROVEMENT.md} \\
\ref{ch:where} & \texttt{rho\_breakdown/\allowbreak{}eval\_protocol\_2026-09-11/}, \texttt{docs/}, \texttt{results/stage0/} & \texttt{02\_PROTOCOL.md}, \texttt{06\_CODE\_FIXES.md}, \texttt{docs/PRIOR\_ART.md} \\
\ref{ch:frame} & \texttt{results/frame/}, \texttt{results/ledger/} & \texttt{docs/PROTOCOL\_RUN.md}, \texttt{TABLE0\_FUNNEL.md} \\
\ref{ch:suppliers} & \texttt{results/elicit/}, \texttt{experiments/remote/} & \texttt{PREREGISTRATION\_SCALE.md} \\
\ref{ch:table4} & \texttt{results/ledger/}, \texttt{src/\allowbreak{}bkrobust/\allowbreak{}benchmarks/\allowbreak{}audit.py} & \texttt{TABLE4.md}, \texttt{CHECKS.md} \\
\ref{ch:pest} & \texttt{results/pest/}, \texttt{results/substrate/} & \texttt{TABLE7.md} \\
\ref{ch:gamma}, \ref{ch:m1b} & \texttt{results/interval/} & \texttt{TABLE1\_L95.md}, \texttt{PREREGISTRATION.md} \\
\ref{ch:tables} & \texttt{rho\_breakdown/\allowbreak{}panelB\_T3\_2026-09-16/}, \texttt{table1\_dossier\_2026-09-16/} & \texttt{PANELB-T3.pdf}, \texttt{TABLE1-DOSSIER.pdf} \\
\ref{ch:review} & \texttt{rho\_breakdown/\allowbreak{}theory\_checks\_2026-09-16/} & the enumeration script \\
\bottomrule
\end{tabular}
\end{center}

Five notes in the dated folders are in Portuguese, because I wrote them for myself on the day.
\texttt{ZE\_EXPERIMENTS.md} summarises each in a paragraph; their content is also in
Chapters~\ref{ch:locality}, \ref{ch:redefinition} and~\ref{ch:where}.

\begin{armadilha}[Absolute paths in the dated folders]
The scripts under \texttt{rho\_breakdown/} were run where they were written and still carry absolute
paths of my machine (for example \texttt{/\allowbreak{}Users/\allowbreak{}josecosta/\allowbreak{}bkrobust/\allowbreak{}src} added to \texttt{sys.path}).
They will not run from a fresh clone without editing that line. \texttt{ZE\_EXPERIMENTS.md} gives the
mapping from each such path to its folder on the branch. The protocol-run scripts in
\texttt{experiments/} do not have this problem: they resolve paths from the repository root.
\end{armadilha}

\prompt[conceito=tabelas-pagas,dificuldade=0.75]{You want to check that the panel-A cells of
Table~1 are what the pipeline produces. Which script, and what does it do when a cell disagrees?}{%
\texttt{rho\_breakdown/\allowbreak{}panelB\_T3\_2026-09-16/\allowbreak{}make\_figs.py}. It reimplements the pipeline from
\texttt{results/\allowbreak{}interval/\allowbreak{}replicates.csv.gz} and \texttt{rows.csv} and asserts, on every run, that it
reproduces the 27 panel-A cells of \texttt{results/\allowbreak{}interval/\allowbreak{}TABLE1\_L95.md}; a disagreement stops the
script.}

\prompt[conceito=gac-backdoor,dificuldade=0.7]{Which file must an audit of adjustment-set validity
call, and which function in the package must it not call?}{It must call the generalised adjustment
criterion in \texttt{src/\allowbreak{}bkrobust/\allowbreak{}benchmarks/\allowbreak{}audit.py}. It must not use
\texttt{is\_valid\_adjustment\_set\_dag}, which implements Pearl's back-door criterion and rejects some
sets that are valid under the generalised criterion. \texttt{tests/\allowbreak{}benchmarks/\allowbreak{}test\_audit\_oracle.py}
guards this.}

\section{Caveats I know about}

\begin{antes}
\antesq[quadro-observavel]{Which stage-0 file quotes a number its own data files no longer give?}
\antesq[screen-leave-k-out]{Which printed table has a selection rule that favours its own headline?}
\end{antes}

Four things about the files deserve a warning before anyone quotes from them.

\begin{medido}[Stage-0 calibration, \texttt{results/\allowbreak{}stage0/\allowbreak{}calibration\_summary.json}]
The corrupted-knowledge calibration on the committed rows was rerun on 13 September after a seeding
fix: \textbf{8,911 rows scored of 16,496 draws}. \texttt{results/\allowbreak{}stage0/\allowbreak{}STAGE0.md} and
\texttt{PREREGISTRATION.md} still quote the first run, \textbf{8,553 rows}, which seeded with Python's
salted \texttt{hash()} and cannot be reproduced across processes. The calibration files on the branch
are the rerun. It changes one reading: on the corrupted rows the hop radius is dangerous on 3 rows
(was 0), the free rule on 9 (was 8), the claim radius on 0 (unchanged).
\end{medido}

Salted \texttt{hash()} seeded six things: the chance arm of the ledger sweep, the Lever~0
differential, the calibration, the first frame-status run, the questionnaire builder and the
source-overlap probe. \texttt{docs/PROTOCOL\_RUN.md} heads its list «Four experiment scripts» and
names five; Chapter~\ref{ch:frame} has the full account. All now seed with SHA-256 of a stable
name. The frozen questionnaire does not regenerate byte for byte
from its builder, so the committed file, with its hash, stays the authoritative input and the
builder refuses to overwrite it.

The second caveat is the selection rule of T3 in \texttt{experiments/\allowbreak{}interval\_tables.py}, analysed
in Chapter~\ref{ch:tables}: it sorts a finite null radius first before taking five rows per network.
The fix is one sort key, or printing all 37 eligible rows. The third is the validity oracle above,
which is still called in two places I have not changed. \texttt{experiments/\allowbreak{}declared\_queries.py}
scores validity at the truth with the back-door function; the committed set judged at the truth is
not the truth's optimal set, so the two criteria can disagree there (they did on 4 of 1,169 ledger
rows). Rerun it with the GAC audit before quoting its validity column. And the hop radius, through
\texttt{hybrid.py}, tests validity with \texttt{mpdag\_criterion.is\_valid\_mpdag}, while the claim
radius uses the generalised criterion; the stricter test can only make $\rhop$ smaller, so the hop
unit's over-certification counts in Chapter~\ref{ch:table4} are, if anything, conservative.
The fourth is what is deliberately absent: \texttt{figures/*} and
\texttt{results/\allowbreak{}axisa3/\allowbreak{}networks/\allowbreak{}acquisition\_manifest.json} were regenerated on my machine by a
reproduction run and are not on the branch, and neither is the pgmpy source tarball.

\prompt[conceito=screen-leave-k-out,dificuldade=0.75]{State the one-line fix for the T3 selection
rule and the alternative that avoids choosing at all.}{In the script that builds T3
(\texttt{interval\_tables.py}, under \texttt{experiments}), remove the term that sorts a finite null
radius first, so the five rows per network are not chosen by the outcome. The alternative: print all
37 eligible rows on the four networks with fitted coefficients instead of choosing.}

\prompt[conceito=pre-registro,dificuldade=0.7]{Why does the committed questionnaire stay
authoritative even though its builder is now reproducible?}{Because the committed questionnaire was
the input every supplier answered and it was hashed before any model saw it. The old builder seeded
with salted \texttt{hash()}, so the file does not regenerate byte for byte; rebuilding it would
produce a different questionnaire from the one the answers refer to. The builder therefore refuses
to overwrite it.}

\naopasse[pre-registro]{Pick one number from Table~4 in Chapter~\ref{ch:table4}. Write down the file
you would open to check it and the script that wrote that file. If you cannot name both, reread the
second table of this chapter.}

\section{What is not here}

The paper drafts are not on the branch: your draft of 16 September is yours, and my outline of
13 September is a planning document I can send separately. Two third-party PDFs that one August
experiment read (Guo and Perković, arXiv 2010.08611; Taeb, Guo and Henckel, arXiv 2511.10625) are
left out; fetch them from arXiv. The response cache of the local suppliers is regenerable from the
questionnaire and stays untracked, as your \texttt{.gitignore} already says.

\begin{exercicio}[metodo=table-reading,conceito=tabelas-pagas]
A reviewer asks where the claim «the certificate counted in claims is dangerous on no deployment
row» comes from. Using the second table of this chapter, name the chapter, the folder, and the two
files you would send, and say which of the two was written first.
\end{exercicio}
\begin{solucao}
Chapter~\ref{ch:table4}; folder \texttt{results/ledger/}; the files \texttt{TABLE4.md} (the
table with the dangerous-row counts) and \texttt{PREREGISTRATION.md} (the predictions and the bucket
definitions). The pre-registration was written first, before the rows were read; the table was
built after the sweep was frozen.
\end{solucao}

\begin{exercicio}[metodo=estimation,conceito=quadro-observavel]
Before merging, you want a rough idea of the conflict surface. Commit~1 touches eight existing files
of the package and adds five. Without running git, estimate which of your recent branches is most
likely to conflict, and what you would check first.
\end{exercicio}
\begin{solucao}
The conflicts can only appear in the eight modified files: \texttt{dagitty.py}, \texttt{measure.py}
and \texttt{parsed.py} under \texttt{benchmarks}, \texttt{conventions.py} under \texttt{core},
\texttt{hybrid.py}, the \texttt{\_\_init\_\_.py} of \texttt{knowledge}, and \texttt{exact.py} and
\texttt{exact\_fast.py} under \texttt{search}. The five added files cannot conflict. A
branch of yours that changed the knowledge-selection or the search dispatch (for example work on
\texttt{measure.py} or \texttt{hybrid.py} after 11 September) is the most likely conflict. Check
first with \texttt{git merge-tree} or a trial merge on a throwaway branch, then look at
\texttt{measure.py}, where both the thinning rule and the rejection labels live.
\end{solucao}

\section*{Chapter map}

\begin{center}
\begin{tikzpicture}[node distance=7mm and 9mm]
\node[cmapc] (br) {branch \texttt{ze/experiments}};
\node[cmap, above left=of br] (aug) {August package\\ \texttt{rho\_breakdown/}};
\node[cmap, above right=of br] (run) {protocol run\\ \texttt{src}, \texttt{experiments}, \texttt{results}};
\node[cmap, below left=of br] (dated) {dated folders\\ 14 Aug to 16 Sep};
\node[cmapg, below right=of br] (idx) {\texttt{ZE\_EXPERIMENTS.md}\\ and this document};
\node[cmapr, below=14mm of br] (cav) {caveats: stage-0 rerun, T3 sort, back-door oracle, absolute paths};
\draw[cedge] (aug) -- node[above, sloped] {merged} (br);
\draw[cedge] (run) -- node[above, sloped] {commits 1--3} (br);
\draw[cedge] (dated) -- node[above, sloped] {commit 4} (br);
\draw[cedge] (idx) -- node[above, sloped] {commit 5} (br);
\draw[cedge] (idx) -- node[right] {lists} (cav);
\draw[cedge] (run) to[bend left=40] node[right, pos=0.7] {checked by 39 tests} (idx);
\end{tikzpicture}
\end{center}
